\section{Residual divisor collisions and private depth}
\label{rc-section}
The results in this section use exact integer identities and finite lattice
geometry. They control prime-power depth repeated across small residuals;
the maximal depth belonging to one residual remains a separate problem.
No logarithmic-form estimate or hypothesis that large primes are simple is
used. Write $\mathcal O=\mathbb Z[\zeta]$, $\zeta^2-\zeta+1=0$,
$P(a+b\zeta)=ab(a+b)$ and $N(a+b\zeta)=a^2+ab+b^2$.

\subsection{Exact cancellation of a common large factor}
Recall the established primitive boundary identities
\begin{equation}\label{rc-old-gcd}
 \gcd(N(x),|P(x)|)=1,\qquad
 \gcd(|P(x)|,|P(yx)|)=\gcd(|P(x)|,|P(y)|),
\end{equation}
where $x$ is primitive and $y$ is arbitrary. For completeness, if $x=(a,b)$
and $y=(r,s)$, the transformed boundary is congruent to
$-P(y)b^3$, $P(y)a^3$, and $P(y)b^3$ modulo $a$, $b$, and $a+b$,
respectively. These three boundary factors are pairwise coprime, so the
unit factors can be cancelled in their gcds. Multiplication of the three
gcd equalities proves the second identity. Reducing the norm modulo each
boundary factor proves the first.

\begin{theorem}[Exact cross-residual overlap]\label{rc-cross}
Let $H,v_1,v_2\in\mathcal O$, suppose $z_i=v_iH$ are primitive, and put
$A_i=|P(z_i)|>0$, $V_i=N(v_i)$. Then
\[
 \gcd(A_1,A_2)=\gcd\bigl(A_1,|P(v_2\bar v_1)|\bigr).
\]
This equality retains all prime-power multiplicities, including those at
two and three.
\end{theorem}
\begin{proof}
In \eqref{rc-old-gcd} take $x=z_1$, $y=v_2\bar v_1$.
Then $yx=V_1z_2$ and $|P(yx)|=V_1^3A_2$.
Since $N(z_1)=V_1N(H)$, the first identity in \eqref{rc-old-gcd}
gives $\gcd(A_1,V_1)=1$. Cancel $V_1^3$ in the left gcd.
\end{proof}

\begin{lemma}[Small residual overlap]\label{rc-height}
If the residuals in Theorem~\ref{rc-cross} are not associates under the six
units, then $P(v_2\bar v_1)\ne0$ and
\[
 \gcd(A_1,A_2)\le(V_1V_2)^{3/2}.
\]
\end{lemma}
\begin{proof}
Each $v_i$ is primitive because $v_iH$ is primitive. A nonzero Eisenstein
integer has zero boundary exactly when it is an integer multiple of a unit.
Thus zero cross-boundary would imply $v_2=(m/V_1)uv_1$ for an integer $m$
and a unit $u$. A rational scalar between primitive integer coordinate pairs
is $1$ or $-1$, by cancelling its denominator and then its numerator.
This would make the residuals associates. Finally the exact cubic identity gives
\[
 3\sqrt3\,|P(D)|=|D^3-\bar D^3|\le2N(D)^{3/2}.
\]
Apply this to $D=v_2\bar v_1$, of norm $V_1V_2$, and use
Theorem~\ref{rc-cross}.
\end{proof}

\subsection{A total budget and exact signed compensation}
Let $v_1,\ldots,v_K$ be pairwise nonassociate residuals, $K\ge1$, with all
$z_i=v_iH$ primitive and $A_i=|P(z_i)|>0$. Set $V_i=N(v_i)$ and
$L=\operatorname{lcm}(A_1,\ldots,A_K)$.

\begin{theorem}[Repeated-multiplicity budget]\label{rc-duplicate}
The integer $D=\prod_iA_i/L$ satisfies
\[
 D\mid\prod_{i<j}\gcd(A_i,A_j),\qquad
 \log D\le\frac32(K-1)\sum_i\log V_i.
\]
In particular $V_i\le B$ implies $\log D\le\frac32K(K-1)\log B$.
\end{theorem}
\begin{proof}
At a fixed prime, arrange its valuations increasingly as $e_1\le\cdots\le e_K$.
The valuation in $D$ is $\sum_{i<K}e_i$, while that in the product of pairwise
gcds is $\sum_{i<K}(K-i)e_i$. This proves the divisibility. Bound each gcd
by Lemma~\ref{rc-height}; each $\log V_i$ occurs in $K-1$ pairs.
\end{proof}

For $Y\ge1$ put $W_Y(A)=\prod_{p\mid A,\ p>Y}p^{v_p(A)-3}$.
Assign each prime dividing $L$ to an index of maximal valuation, breaking ties
by the smallest index. This index is its \emph{owner}; ownership is defined
from actual integer valuations.

\begin{theorem}[Signed compensation]\label{rc-signed}
Define
\[
 D_Y=\prod_{p>Y}p^{\sum_i v_p(A_i)-v_p(L)},\qquad
 E_Y=\prod_{p\mid L,\ p>Y}p^{\#\{i:p\mid A_i\}-1}.
\]
Then $D_Y\mid D$ and
\[
 \prod_iW_Y(A_i)=W_Y(L)\frac{D_Y}{E_Y^3}.
\]
For each index define
\[
 U_i(Y)=\prod_{\substack{p>Y\\i\text{ owns }p}}p^{v_p(A_i)-3},\qquad
 R_i(Y)=\sum_{\substack{p>Y, p\mid A_i\\i\text{ does not own }p}}
                       v_p(A_i)\log p.
\]
Then
\[
 \log W_Y(A_i)\le\log U_i(Y)+R_i(Y),\qquad
 \sum_iR_i(Y)=\log D_Y.
\]
\end{theorem}
\begin{proof}
At a prime occurring in $s$ of the $A_i$, the signed exponent identity is
\[
 \sum_i e_i-3s=(\max_i e_i-3)+(\sum_i e_i-\max_i e_i)-3(s-1).
\]
This proves the product formula. The individual inequality drops only the
negative $-3$ credits at nonowned primes, and their full multiplicities sum
to the exponent of $D_Y$.
\end{proof}

Suppose additionally $H=w^g$, $Q=N(w)\ge7$, and $V_i\le B$.
After rotating each product to the positive sector, let $t_i=\log c_i$;
then $t_i\ge g\log Q/2$. Consequently
\begin{equation}\label{rc-average}
 \frac1K\sum_i\frac{R_i(Y)}{t_i}
 \le\delta:=\frac{3(K-1)\log B}{g\log Q}.
\end{equation}
Apply Markov to the nonnegative quantities $R_i(1)/t_i$. If $\delta>0$,
outside at most a $\sqrt\delta$ proportion of indices, \emph{every} $Y\ge1$
satisfies $R_i(Y)/t_i\le\sqrt\delta$. If $\delta=0$, every contribution is zero.
Thus $K\log B=o(g\log Q)$ gives a simultaneous sublinear repeated-depth
cost for most residuals. No Markov inequality is applied to a signed aggregate:
negative credit could conceal positive individual signed costs.

\begin{corollary}[A specified packet]\label{rc-packet}
If $V_i\le B$ and an integer $q>B^3$ divides two of the $A_i$, those
residuals are associates. In a pairwise nonassociate family, at most one member
can therefore contain a specified such packet.
\end{corollary}
\begin{proof}
Otherwise $q\le\gcd(A_i,A_j)\le(V_iV_j)^{3/2}\le B^3$.
\end{proof}

\subsection{Annotated congruences and a lattice of exact index}
Let $q=\prod p^{e_p}$ be coprime to $N(H)$. At each prime choose one of the
three boundary linear forms $\ell_{\sigma(p)}\in\{a,b,a+b\}$ and define
\[
 \mathcal L(q,\sigma;H)=
 \{v\in\mathbb Z^2:\ell_{\sigma(p)}(vH)\equiv0\pmod{p^{e_p}}
                              \text{ for all }p\mid q\}.
\]

\begin{theorem}[Annotated lattice]\label{rc-lattice}
This lattice has index $q$. Any two rationally independent $v,v'$ in it satisfy
\[
 N(v)N(v')\ge\frac34q^2.
\]
If $q>2B/\sqrt3$, all primitive vectors in it of norm at most $B$ belong
to at most one opposite pair.
\end{theorem}
\noindent\textit{Proof.}\quad For $H=r+s\zeta$, multiplication has matrix
\[
 \begin{pmatrix}r&-s\\s&r+s\end{pmatrix}
\]
of determinant $N(H)$, invertible modulo each retained prime power.
Each chosen row is unimodular modulo that prime, so its map onto
$\mathbb Z/p^{e_p}\mathbb Z$ is surjective. The Chinese remainder theorem
on the two coordinates makes the combined map surjective, proving the index.
The same row vanishes on both vectors modulo each prime power, so
$q\mid\det(v,v')$. Write $v=(a,b)$, $v'=(a',b')$. The exact identity
\[
 4N(v)N(v')-(2aa'+ab'+ba'+2bb')^2=3(ab'-ba')^2
\]
gives the norm-product bound. Primitive integer vectors on one rational
line differ only by sign.\qed\par

\begin{corollary}[Fixed-packet capacity]\label{rc-capacity}
If $\gcd(q,N(H))=1$ and $q>2B/\sqrt3$, at most $2\cdot3^{\omega(q)}$
primitive integer pairs of norm at most $B$ satisfy $q\mid P(vH)$.
For $B\ge1$ this is at most $3^{\omega(q)-1}$ associate classes.
\end{corollary}
\begin{proof}
At each prime, primitive coordinates and invertible multiplication imply that
exactly one boundary arm contains the full prime power. There are at most
$3^{\omega(q)}$ annotations, each with at most one opposite pair by
Theorem~\ref{rc-lattice}. The raw solution set is freely closed under the
six units, which preserve norm and boundary divisibility. Divide its bound
by six. The threshold implies $q>1$ when $B\ge1$.
\end{proof}
The packet in this corollary is specified. A sum over packets chosen after
seeing the residual needs an additional bound on that collection. Likewise,
the lattice theorem constrains the product of two independent solution heights;
it does not lower-bound the first minimum. A lattice may have one very short
primitive vector and a very long independent vector.

\subsection{A pointwise repeated-depth bound}
\begin{theorem}[Linear prime-power collision]\label{rc-linear}
For nonassociate residuals in Theorem~\ref{rc-cross}, any common prime satisfies
\[
 p^{\min(v_p(A_1),v_p(A_2))}\le\frac2{\sqrt3}\sqrt{V_1V_2}.
\]
\end{theorem}
\begin{proof}
A common boundary prime divides neither residual norm. Thus the cross element
$D=v_2\bar v_1$ has unit norm modulo that prime, and at most one of its three
boundary factors is divisible by it. The whole common prime power in
Theorem~\ref{rc-cross} lies in that one nonzero factor.
Each factor has absolute value at most $2\sqrt{N(D)}/\sqrt3$, by
\[
 4N(a,b)=3a^2+(a+2b)^2=3b^2+(2a+b)^2
       =3(a+b)^2+(a-b)^2.
\]
Use $N(D)=V_1V_2$.
\end{proof}

\begin{theorem}[Pointwise nonowner budget]\label{rc-pointwise}
When $V_i\le B$, every index and every $Y\ge1$ satisfy
\[
 R_i(Y)\le\log\operatorname{lcm}(1,\ldots,\lfloor2B/\sqrt3\rfloor)<4B.
\]
For $H=w^g$ as above, $R_i(Y)/t_i\le8B/(g\log Q)$. Hence
$B=o(g\log Q)$ gives sublinear nonowner cost at every residual, with no
exceptional indices.
\end{theorem}
\begin{proof}
For a nonowned prime the owner has at least the same depth, so
Theorem~\ref{rc-linear} gives $p^{v_p(A_i)}\le X=2B/\sqrt3$.
Their product divides $\operatorname{lcm}(1,\ldots,\lfloor X\rfloor)$,
whose logarithm is $\psi(X)=\sum_{p^e\le X}\log p$.
Every prime power in $(m,2m]$ contributes to the corresponding valuation of
$\binom{2m}{m}$, with nonnegative contributions at other levels. Thus
$\psi(2m)-\psi(m)\le2m\log2$. Dyadic summation yields
$\psi(X)<4X\log2\le3X$, where $\log2\le3/4$ follows from the trapezoidal
upper bound for $\int_1^2dx/x$. Therefore $\psi(X)\le2\sqrt3B<4B$.
The height lower bound proves the final assertion.
\end{proof}

\subsection{An exact limit of the collision argument}
Take $w=2+\zeta$, $v_1=1$, $v_2=1+5\zeta$ and $g_k=2\cdot5^k$ for
$k\ge1$. The products are primitive since the residual norm $31$ is disjoint
from $N(w)=7$. Their displayed lambda parameters tend to zero. The elementary
homogeneous rank law at five has rank two and first depth one:
$w^2=3+5\zeta$ has boundary $120$. Hence
\[
 v_5(|P(w^{g_k})|)=k+1.
\]
But $P(v_2)=30$, so Theorem~\ref{rc-cross} forces
$v_5(|P(v_2w^{g_k})|)=1$. The maximal depth is unbounded while its repeated
depth is exactly one. This disproves only the attempted inference that a
collision budget bounds maximal depth. The fixed prime five is eventually
below the moving lambda cutoff, so the family does not refute that tail gate.

The exact signed formulas, capacity bound and pointwise nonowner theorem leave
the owner packets $U_i(Y)$ unbounded. For primes greater than $2B/\sqrt3$,
all supported depth is already private to one associate class. The new
arithmetic task is to constrain that private first short direction using the
special common factor $w^g$, without replacing a specified packet by an
uncontrolled union or assuming prime simplicity. The companion finite residue
measure section develops one such reference object and states its transfer
gap explicitly. The results here do not establish standard $abc$.
